\documentclass[10pt]{article}
\usepackage[T1]{fontenc}
\usepackage[utf8]{inputenc}
\usepackage[french]{babel}
\usepackage{amsmath,amssymb,mathrsfs,mathpazo}
\usepackage{geometry}
\geometry{a4paper,margin=2cm}
\usepackage{graphicx,array,tabularx,tabularray,colortbl,multirow,multicol,booktabs,enumitem,pifont,ulem,fancybox,diagbox,cancel}
\usepackage{pstricks,pstricks-add,pst-plot,pst-tree,pst-node,pst-text,pst-func,pst-math,pst-bspline,pst-3dplot,pst-eucl,pst-solides3d}
\let\pscalculate\undefined
\usepackage{pst-bezier}
\usepackage{tikz,pgfplots,tkz-tab}
\usetikzlibrary{arrows,arrows.meta,patterns,calc,positioning,trees,decorations.markings,shapes.geometric,shapes.misc,shapes.arrows,intersections,decorations.pathmorphing,decorations.pathreplacing,backgrounds,shadows,babel}
\pgfplotsset{compat=1.18}
\usepgfplotslibrary{fillbetween}
\usepackage[np]{numprint}
\usepackage{listings,makecell,pas-tableur,twoopt}
\let\barre\undefined
\usepackage[pstricks]{tablvar}
\definecolor{darkgreen}{RGB}{0,100,0}
\providecommand{\red}{\color{red}}
\providecommand{\blue}{\color{blue}}
\providecommand{\green}{\color{green}}
\providecommand{\black}{\color{black}}
\providecommand{\white}{\color{white}}
\providecommand{\gray}{\color{gray}}
\providecommand{\cyan}{\color{cyan}}
\providecommand{\magenta}{\color{magenta}}
\providecommand{\yellow}{\color{yellow}}
\providecommand{\violet}{\color{violet}}
\providecommand{\orange}{\color{orange}}
\providecommand{\euro}{\texteuro}
\providecommand{\EUR}{\texteuro}
\providecommand{\ieme}{\textsuperscript{e}}
\providecommand{\degres}{\ensuremath{^\circ}}
\providecommand{\degree}{\ensuremath{^\circ}}
\providecommand{\vv}[1]{\overrightarrow{#1}}
\providecommand{\Coord}[1]{#1}
\providecommand{\decosix}{}
\providecommand{\starredbullet}{\ensuremath{\bullet}}
\pagestyle{empty}
\setlength{\parindent}{0pt}
\usepackage{adjustbox}\DeclareUnicodeCharacter{2236}{:}\DeclareUnicodeCharacter{2190}{\ensuremath{\leftarrow}}
\providecommand{\R}{\mathbb{R}}
\providecommand{\N}{\mathbb{N}}
\providecommand{\Z}{\mathbb{Z}}
\providecommand{\Q}{\mathbb{Q}}
\providecommand{\D}{\mathbb{D}}
\providecommand{\C}{\mathbb{C}}
\providecommand{\e}{{\mathrm{e}}}
\providecommand{\vect}[1]{\overrightarrow{#1}}
\providecommand{\vectt}[1]{\overrightarrow{\mathrm{#1}}}
\providecommand{\vvt}[1]{\overrightarrow{\mathrm{#1}}}
\providecommand{\np}[2][]{{#2~#1}}
\providecommand{\vv}[1]{{\overrightarrow{#1}}}
\renewcommand{\d}{{\mathrm{d}}}
\providecommand{\degres}{{^\circ}}
\providecommand{\degre}{{^\circ}}
\providecommand{\blue}{\color{blue}}
\providecommand{\red}{\color{red}}
\providecommand{\enskip}{\quad}
\providecommand{\llbracket}{\lbrack\!\lbrack}
\providecommand{\rrbracket}{\rbrack\!\rbrack}
\providecommand{\euro}{{\text{€}}}
\providecommand{\no}{{\text{n°}}}
\providecommand{\interval}[2]{{[#1;#2]}}
\providecommand{\intervFO}[2]{{[#1;#2[}}
\providecommand{\dotfill}{\ldots}

\providecommand{\ds}{\displaystyle}
\providecommand{\textding}[1]{#1}
\providecommand{\Oij}{\(\mathrm{O};\vec{\imath},\vec{\jmath}\)}
\providecommand{\Oijk}{\(\mathrm{O};\vec{\imath},\vec{\jmath},\vec{k}\)}
\providecommand{\Ouv}{\(\mathrm{O};\vec{u},\vec{v}\)}
\providecommand{\intervFF}[2]{\ensuremath{[#1;#2]}}
\renewcommand{\euro}{\ifmmode\text{\texteuro}\else\texteuro\fi}
\DeclareUnicodeCharacter{00B0}{\ifmmode^{\circ}\else\textdegree\fi}
\DeclareUnicodeCharacter{20AC}{\ifmmode\text{\texteuro}\else\texteuro\fi}
% Énoncés : Éducation nationale. Transcription et corrigés : APMEP, Denis Vergès et contributeurs. Découpage et classement : CoopMaths, Sébastien Lozano, Éric Elter et contributeurs. Adaptation de la présentation : Maths Explorer.
% CC BY-SA : https://coopmaths.fr/www/about/
% https://coopmaths.fr/alea/static/bac/2025/tex/bac_2025_09_sujet2_metropole_3.tex
\begin{document}
\noindent\textit{APMEP / CoopMaths. CC BY-SA. Présentation adaptée par Maths Explorer.}\par\medskip
\medskip

On considère la fonction $f$ définie sur ]0~;~8] par 

\[f(x) = \dfrac{10\ln \left(- x^2 + 7x + 9\right)}{x}\]

Soit $C_f$ la représentation graphique de la fonction $f$ dans un repère orthonormé \Oij.

\medskip

\textbf{Partie A}

\medskip

\begin{enumerate}
\item Résoudre dans $\R$ l’inéquation $- x^2 + 7x + 8 \geqslant 0$.
\item En déduire que pour tout $x \in ]0~;~8]$, on a $f(x) \geqslant 0$.
\item Interpréter graphiquement ce résultat.
\end{enumerate}

\medskip

\textbf{Partie B}

La courbe $C_f$ est représentée ci-dessous.

Soit $M$ le point de $C_f$ d’abscisse $x$ avec $x \in ]0~;~8]$.

On appelle $N$ et $P$ les projetés orthogonaux du point $M$ respectivement sur l’axe des abscisses et sur l’axe des ordonnées.

Dans cette partie, on s’intéresse à l’aire $\mathcal{A}(x)$ du rectangle O$NMP$.

\begin{center}
\psset{unit=1cm,arrowsize=2pt 3}
\begin{pspicture}(-0.5,-0.5)(8.5,10)
\uput[ur](4,7.6){\red $(C_f)$}
\psaxes[linewidth=1.25pt,labelFontSize=\scriptstyle,labels=none](0,0)(8.5,10)
\psaxes[linewidth=1.25pt,labelFontSize=\scriptstyle,labels=none]{->}(1,1)
\psframe[linewidth=1.5pt](6.8,3.438)\psframe(6.8,0)(6.6,0.2)\psframe(0,3.438)(0.2,3.238)
\uput[dl](0,0){O}\uput[l](0,3.438){$P$}\uput[ur](6.8,3.438){$M$}\uput[d](6.8,0){$N$}
\psgrid[subgriddiv=1,gridwidth=0.3pt,labels=none](0,0)(8,10)
\psplot[plotpoints=2000,linewidth=1.25pt,linecolor=red]{3}{8}{7 x mul 9 add x dup mul sub ln 10 mul x div}
\uput[d](0.5,0){\small $\vect{\imath}$}\uput[l](0,0.5){\small $\vect{\jmath}$}
\end{pspicture}
\end{center}
\medskip

\begin{enumerate}
\item Donner les coordonnées des points $N$ et $P$ en fonction de $x$.
\item Montrer que pour tout $x$ appartenant à l’intervalle ]0~;~8],

\[\mathcal{A}(x) = 10 \ln \left(-x^2 + 7x + 9\right)\]

\item Existe-t-il une position du point $M$ pour laquelle l’aire du rectangle O$NMP$ est
maximale ? Si elle existe, déterminer cette position.
\end{enumerate}

\medskip

\textbf{Partie C}

\medskip

On considère un réel strictement positif $k$.

On souhaite déterminer la plus petite valeur de $x$, approchée au dixième, appartenant à [3,5~;~8] pour laquelle l’aire $\mathcal{A}(x)$ devient inférieure ou égale à $k$.

Pour ce faire, on considère l’algorithme ci-dessous.

Pour rappel, en langage Python, $\ln (x)$ s’écrit log ($x$).

\begin{center}
\begin{ttfamily}
\begin{tabular}{c l}
1&from math import *\\
2&\\
3&def A(x) :\\
4&\quad return 10*log (- 1* x**2 + 7*x + 9)\\
5&\quad \\
6&def pluspetitevaleur(k) :\\
7&\quad x = 3.5\\
8&\quad while A(x).......... :\\
9& \qquad x = x + 0.1\\
10&\quad return ...........
\end{tabular}
\end{ttfamily}
\end{center}

\smallskip

\begin{enumerate}
\item Recopier et compléter les lignes 8 et 10 de l’algorithme.
\item Quel nombre renvoie alors l’instruction \texttt{pluspetitevaleur}(30) ?
\item Que se passe-t-il lorsque $k = 35$ ? Justifier.
\end{enumerate}

\medskip

\end{document}
